
\documentclass{article}
\usepackage{colortbl}
\usepackage{makecell}
\usepackage{multirow}
\usepackage{supertabular}

\begin{document}

\newcounter{utterance}

\centering \large Interaction Transcript for game `hot\_air\_balloon', experiment `air\_balloon\_survival\_en\_reasoning off\_hard', episode 4 with reka{-}flash{-}3.1{-}t0.0.
\vspace{24pt}

{ \footnotesize  \setcounter{utterance}{1}
\setlength{\tabcolsep}{0pt}
\begin{supertabular}{c@{$\;$}|p{.15\linewidth}@{}p{.15\linewidth}p{.15\linewidth}p{.15\linewidth}p{.15\linewidth}p{.15\linewidth}}
    \# & $\;$A & \multicolumn{4}{c}{Game Master} & $\;\:$B\\
    \hline

    \theutterance \stepcounter{utterance}  
    & & \multicolumn{4}{p{0.6\linewidth}}{
        \cellcolor[rgb]{0.9,0.9,0.9}{
            \makecell[{{p{\linewidth}}}]{
                \texttt{\tiny{[P1$\langle$GM]}}
                \texttt{You are participating in a collaborative negotiation game.} \\
\\ 
\texttt{Together with another participant, you must agree on a single set of items that will be kept. Each of you has your own view of how much each item matters to you (importance). You do not know how the other participant values the items. Additionally, you are given the effort each item demands.} \\
\texttt{You may only agree on a set if the total effort of the selected items does not exceed a shared limit:} \\
\\ 
\texttt{LIMIT = 2984} \\
\\ 
\texttt{Here are the individual item effort values:} \\
\\ 
\texttt{Item effort = \{"A21": 350, "B93": 276, "A11": 616, "B33": 737, "C08": 531, "B39": 389, "B96": 24, "C85": 125, "A82": 338, "B31": 356, "B20": 143, "A60": 117, "B09": 257, "C81": 921, "B21": 789\}} \\
\\ 
\texttt{Here is your personal view on the importance of each item:} \\
\\ 
\texttt{Item importance values = \{"A21": 138, "B93": 583, "A11": 868, "B33": 822, "C08": 783, "B39": 65, "B96": 262, "C85": 121, "A82": 508, "B31": 780, "B20": 461, "A60": 484, "B09": 668, "C81": 389, "B21": 808\}} \\
\\ 
\texttt{Goal:} \\
\\ 
\texttt{Your goal is to negotiate a shared set of items that benefits you as much as possible (i.e., maximizes total importance to YOU), while staying within the LIMIT. You are not required to make a PROPOSAL in every message {-} you can simply negotiate as well. All tactics are allowed!} \\
\\ 
\texttt{Interaction Protocol:} \\
\\ 
\texttt{You may only use the following structured formats in a message:} \\
\\ 
\texttt{PROPOSAL: \{'A', 'B', 'C', …\}} \\
\texttt{Propose keeping exactly those items.} \\
\texttt{REFUSE: \{'A', 'B', 'C', …\}} \\
\texttt{Explicitly reject opponent's proposal.} \\
\texttt{ARGUMENT: \{'...'\}} \\
\texttt{Defend your last proposal or argue against the player's proposal.} \\
\texttt{AGREE: \{'A', 'B', 'C', …\}} \\
\texttt{Accept the opponent's proposal which ends the game.} \\
\\ 
\\ 
\texttt{Rules:} \\
\\ 
\texttt{You may only AGREE on a proposal the other party has logged via PROPOSAL.} \\
\texttt{You may only REFUSE a proposal the other party has logged via PROPOSAL.} \\
\texttt{Total effort of any PROPOSAL or AGREE set must be ≤ LIMIT.} \\
\texttt{Do NOT reveal your hidden importance scores.} \\
\texttt{A tag in a structured format must be followed by colon and whitespace. The argument must be a python set containing 0 or more strings.} \\
\texttt{So, it must be of the form TAG: \{...\}} \\
\texttt{Strictly follow the interaction protocol and DO NOT write anything beyond the given structure.} \\
\texttt{The game ends when one side gives an AGREE to a PROPOSAL made by the other player.} \\
\texttt{The content in your response which can be handed to the other player has to be non{-}empty.} \\
\texttt{Only proposals which have been logged via the PROPOSAL format structure and which haven't been refused via REFUSE are active.} \\
\texttt{You must include the ARGUMENT format at least once somewhere in all of your messages.} \\
            }
        }
    }
    & \\ \\

    \theutterance \stepcounter{utterance}  
    & \multicolumn{4}{p{0.6\linewidth}}{
        \cellcolor[rgb]{0.9,0.9,0.9}{
            \makecell[{{p{\linewidth}}}]{
                \texttt{\tiny{[P1$\rangle$GM]}}
                \texttt{<reasoning>} \\
\texttt{Let's tackle this negotiation problem. The goal is to maximize my total importance while keeping the total effort under 2984. First, I need to figure out which items I value the most and which have the least effort. Since I don't know the other person's preferences, I should focus on my own importance scores but also consider the effort to stay within the limit.} \\
\\ 
\texttt{Let me list out all the items with their effort and my importance:} \\
\\ 
\texttt{Effort:} \\
\texttt{A21: 350} \\
\texttt{B93: 276} \\
\texttt{A11: 616} \\
\texttt{B33: 737} \\
\texttt{C08: 531} \\
\texttt{B39: 389} \\
\texttt{B96: 24} \\
\texttt{C85: 125} \\
\texttt{A82: 338} \\
\texttt{B31: 356} \\
\texttt{B20: 143} \\
\texttt{A60: 117} \\
\texttt{B09: 257} \\
\texttt{C81: 921} \\
\texttt{B21: 789} \\
\\ 
\texttt{Importance:} \\
\texttt{A21: 138} \\
\texttt{B93: 583} \\
\texttt{A11: 868} \\
\texttt{B33: 822} \\
\texttt{C08: 783} \\
\texttt{B39: 65} \\
\texttt{B96: 262} \\
\texttt{C85: 121} \\
\texttt{A82: 508} \\
\texttt{B31: 780} \\
\texttt{B20: 461} \\
\texttt{A60: 484} \\
\texttt{B09: 668} \\
\texttt{C81: 389} \\
\texttt{B21: 808} \\
\\ 
\texttt{First, I should sort the items based on my importance in descending order to prioritize the ones I value the most. Let's do that:} \\
\\ 
\texttt{1. B33: 822} \\
\texttt{2. B21: 808} \\
\texttt{3. A11: 868 (Wait, wait, A11's importance is 868? Let me check again. Yes, A11 is 868. So that's higher than B33 and B21. Wait, no, B33 is 822, B21 is 808, A11 is 868. So A11 is first. Then B33, then B21. Then C08:783, B09:668, A82:508, B20:461, B96:262, C85:121, B39:65, A21:138, A60:117, C81:389, B93:583. Wait, let me list them properly:} \\
\\ 
\texttt{Sorted by importance (mine) descending:} \\
\\ 
\texttt{A11: 868} \\
\texttt{B33: 822} \\
\texttt{B21: 808} \\
\texttt{C08: 783} \\
\texttt{B} \\
            }
        }
    }
    & & \\ \\

    \theutterance \stepcounter{utterance}  
    & & \multicolumn{4}{p{0.6\linewidth}}{
        \cellcolor[rgb]{0.9,0.9,0.9}{
            \makecell[{{p{\linewidth}}}]{
                \texttt{\tiny{[P1$\langle$GM]}}
                \texttt{Your response did not contain an argument. You must include the structured format ARGUMENT: \{...\} somewhere in your response. Try again} \\
            }
        }
    }
    & \\ \\

    \theutterance \stepcounter{utterance}  
    & \multicolumn{4}{p{0.6\linewidth}}{
        \cellcolor[rgb]{0.9,0.9,0.9}{
            \makecell[{{p{\linewidth}}}]{
                \texttt{\tiny{[P1$\rangle$GM]}}
                \texttt{<reasoning>} \\
\texttt{Let me try to figure out how to approach this. The user mentioned that my previous response didn't include an ARGUMENT. Let me check the requirements again.} \\
\\ 
\texttt{The interaction protocol says I must use structured formats like PROPOSAL, REFUSE, ARGUMENT, AGREE. The ARGUMENT must be a Python set containing 0 or more strings. Also, I need to include ARGUMENT at least once in my messages. } \\
\\ 
\texttt{In my previous attempt, I proposed a set of items and included an ARGUMENT. Wait, maybe I made a formatting error. Let me look at the example. The user's example shows ARGUMENT: \{'...'\}, where the argument is a set. But in my response, I wrote ARGUMENT: \{'High importance items included:...'\}, which is a string inside the set. Wait, no, the user's example shows the argument as a set containing strings. Wait, the instruction says: "The argument must be a python set containing 0 or more strings." So each element in the set should be a string. But in my previous response, I had ARGUMENT: \{'High importance items included:...'\}, which is a single string element. That's allowed, right? Because a set can have one element. But maybe the user expects each element to be a separate point. Wait, the example given in the problem statement shows ARGUMENT: \{'...'\}, but the actual content inside is a string. Wait, looking back:} \\
\\ 
\texttt{The problem says: "The argument must be a python set containing 0 or more strings. So it must be of the form TAG: \{...\}" where the... are strings. So each element in the set is a string. So in my previous response, I had ARGUMENT: \{'High importance items included:...'\}, which is a set with one string element. That should be correct. But maybe the user's system is expecting a different format. Alternatively, perhaps I made a mistake in the syntax. Let me check the exact example from the problem statement.} \\
\\ 
\texttt{In the problem statement, the example shows:} \\
\\ 
\texttt{ARGUMENT: \{'...'\}} \\
\\ 
\texttt{But the actual content inside the set is a string. Wait, no. Wait, the problem says: "the argument must be a python set containing 0 or more strings." So each element in the set is a string. So for example, ARGUMENT: \{'This item is crucial', 'Effort is within limit'\} would be two strings. But in} \\
            }
        }
    }
    & & \\ \\

    \theutterance \stepcounter{utterance}  
    & & \multicolumn{4}{p{0.6\linewidth}}{
        \cellcolor[rgb]{0.9,0.9,0.9}{
            \makecell[{{p{\linewidth}}}]{
                \texttt{\tiny{[P1$\langle$GM]}}
                \texttt{Your response contained an untagged sequence or you used STRATEGIC REASONING more than once. You may only use the structured formats as explained in the initial message. They must all be of the form TAG: \{...\}} \\
            }
        }
    }
    & \\ \\

    \theutterance \stepcounter{utterance}  
    & \multicolumn{4}{p{0.6\linewidth}}{
        \cellcolor[rgb]{0.9,0.9,0.9}{
            \makecell[{{p{\linewidth}}}]{
                \texttt{\tiny{[P1$\rangle$GM]}}
                \texttt{<reasoning>} \\
\texttt{Let me try to figure out what's wrong here. The user mentioned that my response had an untagged sequence or used STRATEGIC REASONING more than once. The error message says I can only use the structured formats like PROPOSAL, REFUSE, ARGUMENT, AGREE, each in the form TAG: \{...\}.} \\
\\ 
\texttt{Looking back at my previous response, I included STRATEGIC REASONING: \{...\} which wasn't one of the allowed tags. The allowed tags are PROPOSAL, REFUSE, ARGUMENT, AGREE. So I mistakenly added a new tag, which is not allowed. That's why the system rejected it. } \\
\\ 
\texttt{I need to make sure that all parts of my message are strictly using the allowed tags. The user also mentioned that I must include ARGUMENT at least once. In my previous attempt, I had ARGUMENT, but maybe I included another tag incorrectly. Let me check again.} \\
\\ 
\texttt{In the response that caused the error, I had:} \\
\\ 
\texttt{PROPOSAL: \{'B21', 'B33', 'A11', 'C08', 'B09', 'B96', 'B20', 'A60', 'B93', 'C85', 'A82', 'B31'\}} \\
\\ 
\texttt{ARGUMENT: \{'High importance items included: B21 (808), B33 (822), A11 (868), C08 (783), B09 (668). Effort total: 2984'\}} \\
\\ 
\texttt{But wait, the user's error message says "You may only use the structured formats as explained in the initial message. They must all be of the form TAG: \{...\}". The allowed tags are PROPOSAL, REFUSE, ARGUMENT, AGREE. So if I included any other tags like STRATEGIC REASONING, that's invalid. Wait, in my example, I didn't use STRATEGIC REASONING. Wait, maybe the user's system is case{-}sensitive? Or perhaps I used a different capitalization. Let me check.} \\
\\ 
\texttt{Wait, in the problem statement, the allowed tags are PROPOSAL, REFUSE, ARGUMENT, AGREE. So the tags must be exactly those, with the correct capitalization. If I used lowercase or different casing, that might be an issue. But in my example, I used uppercase. Hmm.} \\
\\ 
\texttt{Wait, the user's error message says "You may only use the structured formats as explained in the initial message. They must all} \\
            }
        }
    }
    & & \\ \\

    \theutterance \stepcounter{utterance}  
    & & & \multicolumn{2}{p{0.3\linewidth}}{
        \cellcolor[rgb]{0.9,0.9,0.9}{
            \makecell[{{p{\linewidth}}}]{
                \texttt{\tiny{[GM$|$GM]}}
                \texttt{abort} \\
            }
        }
    }
    & & \\ \\

    \theutterance \stepcounter{utterance}  
    & & & \multicolumn{2}{p{0.3\linewidth}}{
        \cellcolor[rgb]{0.9,0.9,0.9}{
            \makecell[{{p{\linewidth}}}]{
                \texttt{\tiny{[GM$|$GM]}}
                \texttt{end game} \\
            }
        }
    }
    & & \\ \\

\end{supertabular}
}

\end{document}
